% A neural surrogate for the sodium-boiling phase of an SFR ULOF transient.
%
%   latexmk -pdf paper/paper.tex        (or: pdflatex paper/paper.tex, twice)
%
% LAYOUT: IOP Journal of Physics: Conference Series, from __DEV/paper_template.docx.
% The template is a Word *layout guide*, not a submission binary, so this file encodes
% its specification the same way `iop-jpcs.typ` does -- every constant below is taken
% from the guide's Table 1 (margins) and Table 2 (section styles), and the two files
% must be changed together. Deliberately NOT built on `iopart.cls`: the class imposes
% its own geometry, which would silently override the guide it is supposed to follow,
% and it is not installed everywhere. If IOP later ask for `iopart.cls` specifically,
% the body below transfers unchanged -- only this preamble is discarded.
%
% Content is kept in step with `paper.md` by hand. Where they disagree,
% `docs/axial_nn.md` holds the measurement record and both are wrong until reconciled
% against it.
\documentclass[11pt,a4paper]{article}

% --- Table 1 of the guide: A4, top 4.0, bottom 2.7, left 2.5, right 2.5 cm, and no
% --- headers, footers or page numbers -- production adds those.
\usepackage[a4paper,top=4.0cm,bottom=2.7cm,left=2.5cm,right=2.5cm]{geometry}
\pagestyle{empty}

% Times at 11 pt. `newtx` if usable, `mathptmx` as the fallback that ships everywhere.
%
% BOTH files are tested, and the second one is the point: `newtxtext.sty` inputs
% `binhex.tex`, which Debian ships in a DIFFERENT package (`texlive-plain-generic`).
% Testing only for the style file finds newtx present, selects it, and then dies on the
% missing input -- which is what happened on this lane's first run. A guard that checks
% the wrong file is not a guard.
\usepackage{ifthen}
\IfFileExists{newtxtext.sty}{%
  \IfFileExists{binhex.tex}{\usepackage{newtxtext,newtxmath}}{\usepackage{mathptmx}}%
}{\usepackage{mathptmx}}
\usepackage[T1]{fontenc}
\usepackage[utf8]{inputenc}
\usepackage{amsmath,amssymb}
\usepackage{graphicx}
% Drawn by `python -m pinn_sfr_transient.axial.charts`, never committed and never
% exported by hand -- CI regenerates them from the physics before pdflatex runs. Draw
% them first for a local build; the directory is gitignored.
\graphicspath{{docs/img/charts/}}
\usepackage{titlesec}
\usepackage[font=small,labelsep=period,singlelinecheck=off]{caption}
\usepackage{booktabs}

\setlength{\parindent}{5mm}   % first-line indent, per the guide
\setlength{\parskip}{0pt}

% --- Table 2: sections 11 pt bold, subsections 11 pt italic, one line space before and
% --- none after. Subsubsections run into the paragraph and end with a full stop.
\titleformat{\section}{\normalsize\bfseries}{\thesection.}{0.5em}{}
\titlespacing*{\section}{0pt}{\baselineskip}{0pt}
\titleformat{\subsection}{\normalsize\itshape}{\thesubsection}{0.5em}{}
\titlespacing*{\subsection}{0pt}{\baselineskip}{0pt}
\titleformat{\subsubsection}[runin]{\normalsize\itshape}{\thesubsubsection}{0.5em}{}[.]
\titlespacing*{\subsubsection}{0pt}{0.5\baselineskip}{0.4em}

% Captions: above tables, below figures, 10 pt. `caption`'s `font=small` gives 10 pt
% against an 11 pt body; the position is a matter of where \caption is placed.
\captionsetup[table]{position=top}
\captionsetup[figure]{position=bottom}

\begin{document}

% ---------------------------------------------------------------- title block ------
% Title 17 pt bold, flush left, unjustified. Authors and abstract indented 25 mm.
\begingroup
\raggedright
\fontsize{17}{20}\selectfont\bfseries
A neural surrogate for the sodium-boiling phase of an SFR unprotected loss-of-flow
transient: what it reproduces, and what it does not
\par
\endgroup

\vspace{1.2em}

\begingroup
\leftskip=25mm \raggedright
P Pigazzini\textsuperscript{1}
\par
\vspace{0.6em}
\fontsize{10}{12}\selectfont
\textsuperscript{1} Independent researcher\\
Corresponding author e-mail: pasquale.pigazzini@gmail.com
\par
\endgroup

\vspace{1.2em}

\begingroup
\leftskip=25mm \rightskip=25mm
\fontsize{10}{12}\selectfont
\noindent\textbf{Abstract.}
We solve the sodium-boiling phase of an unprotected loss-of-flow (ULOF) transient in a
sodium-cooled fast reactor with a physics-informed neural surrogate, and report it as a
reactor-physics result rather than a machine-learning one. One coolant channel is resolved
axially with four material fields---fuel, cladding, film and coolant temperature---and a
sodium void fraction, coupled to six-group point kinetics through a logarithmic Doppler
integral and a void-worth integral whose weight changes sign near the top of the core.
Thermophysics, the saturation-plus-superheat boiling criterion and both feedback laws are
taken from the SAS4A/SASSYS-1 manual; four deviations from it are registered with their
justification, of which two are load-bearing. Against a verified stiff Radau reference the
surrogate reproduces the temperature fields to a relative $L_2$ of $1.6\times10^{-3}$,
places boiling onset within 0.008\,s of the reference's 10.9784\,s, and reproduces 99.4\%
of the peak voided length with a saturation margin of $+67.4$\,K against $+69.2$\,K. Both
of the first two now sit at the reference solver's own resolution, so we report them as met
rather than as measured accuracies. What the surrogate does \emph{not} yet deliver is the
closed reactivity loop: the void-worth integral is a near-cancellation of two large
opposite contributions, and driving the kinetics from the learned fields recovers only
8--16\% of it, while the non-cancelling Doppler integral over the same fields is correct to
1.017. We show by a dual-weighted-residual argument that this is not a sampling deficiency,
quantify the cancellation, and identify the closed-loop void reactivity as the outstanding
physics problem.
\par
\endgroup

\vspace{1.5em}

% ---------------------------------------------------------------------- body -------
\section{Introduction}

In an unprotected loss-of-flow accident the primary pumps coast down while the reactor
protection system fails to scram. Coolant flow decays, the sodium outlet temperature rises,
and if saturation plus the required superheat is reached the coolant boils. In a
sodium-cooled fast reactor the resulting voiding carries a \emph{positive} reactivity
contribution over most of the core height, and it competes with the prompt negative Doppler
feedback. Whether the excursion terminates benignly is decided by the sign and timing of
that competition, which makes the \emph{onset} and \emph{axial extent} of boiling the
quantities of engineering interest, not merely a field norm.

Whole-plant analysis of this sequence is the province of system codes---SAS4A/SASSYS-1
being the reference of record [1]. A differentiable surrogate for the boiling phase is
attractive for uncertainty propagation, for parameter inference from limited
instrumentation, and for embedding a channel model inside an optimisation loop, all of
which need many evaluations and derivatives with respect to inputs.

This paper reports such a surrogate and is explicit about the division of labour. The
machine-learning system is described only as far as needed to reproduce it
(section~4); the detail is spent on the physics being solved (section~2), on the reference
against which it is judged and the resolution that reference actually has (section~3), on
what the surrogate reproduces (section~5), and on the one coupling it does not yet close
(section~6). The complete methodological record---architecture sweeps, optimiser
comparisons, seed statistics, cross-implementation checks and the negative results---is
published alongside in the repository documentation and is cited rather than reproduced
here.

\section{Physical model}

\subsection{Geometry and state}

A single coolant channel is resolved along its axis, $\zeta\in[0,1]$, with radially lumped
materials. The state at each height is four temperatures---fuel $T_f$, cladding $T_{cl}$,
film $T_s$ and coolant $T_c$---together with a void fraction $\alpha$. Radial lumping
(D-GEOM-1) and retention of a lumped structure node (D-GEOM-2) follow the manual's own
reduced treatment.

\subsection{Governing equations}

Fuel and cladding conduction follow manual Eq.~3.3-1, with the fuel-to-cladding flux of
Eq.~3.3-4 carrying \emph{both} gap conductance and radiation. The coolant energy equation
is Eq.~3.3-5 with its three sources $Q_c + Q_{ec} + Q_{sc}$, and includes direct neutron
and gamma heating $\gamma_c$ (Eq.~3.3-6). Coolant transport is advective with a prescribed
flow decay,
\begin{equation}
f(t) = f_{\mathrm{nc}} + (1 - f_{\mathrm{nc}})\,e^{-t/\tau_{\mathrm{pump}}},
\end{equation}
with $f_{\mathrm{nc}}$ the natural-circulation floor. Pre-boiling momentum uses the
manual's Eq.~3.9-1 result that $w = w(t)$ is independent of height.

\subsection{Sodium properties and the boiling criterion}

All thirteen numbered relations of manual section~12.13 are implemented. Their stated
validity is \textbf{590--2270\,K}, the fits stopping near 90\% of the critical point
(2503.3\,K) because $C_l$, $\beta_s$ and $\alpha_p$ all contain $1/(T_c - T)$ and diverge
there. The implementation neither clamps nor raises---a hard guard inside a residual would
abort a differentiable solve the moment a transient overshot---and instead exposes a range
check that the reference solution asserts.

Boiling onset is the saturation-plus-superheat criterion of section~12.4: the coolant boils
where $T_c > T_{\mathrm{sat}}(p) + \Delta T_{\mathrm{sup}}$. The cladding-and-structure to
vapour heat path follows section~12.5.1.

\subsection{Neutronics and feedback}

Power follows point kinetics with six delayed-neutron precursor groups (Eq.~4.3-1) closed
by the prompt-jump approximation (Eq.~4.2-4). Two feedbacks couple the thermal fields back
to the power, both as \emph{axial integrals}: \emph{Doppler}, logarithmic in fuel
temperature and interpolated between flooded and voided states (Eqs.~4.5-2, 4.5-3); and
\emph{coolant density and void as a single worth distribution} (Eq.~4.5-25),
\begin{equation}\label{eq:void}
\rho_{\mathrm{void}}(t) = \int_0^1 w(\zeta)\,\alpha(\zeta,t)\,\mathrm{d}\zeta ,
\end{equation}
where the worth $w$ is \textbf{positive over most of the core and negative near the top}.
That sign change is the physical origin of the difficulty reported in section~6.

\subsection{Registered deviations from the manual}

Every departure from SAS4A is registered with its justification; four matter here.

\subsubsection{D-KIN-1, prompt jump}
The prompt neutron mode is eliminated algebraically. Standard practice, and it removes the
stiffest timescale from the system.

\subsubsection{D-TH-3, the void slaved to temperature}
Rather than tracking vapour dynamics, the void fraction is closed algebraically on the
coolant temperature,
\begin{equation}\label{eq:closure}
\alpha = 1 - \bigl(1 - b(T_c)\bigr)^{3},
\end{equation}
with $b$ the superheat switch. Justified by timescale separation---vapour fills a node in
0.71\,ms against an advective 0.113\,s---and selected from six candidate closures on
maximum voided-length error, at 1.2\% against the reference's own 2.6\% mesh error. The
cubic form is chosen for its finite slope at $b = 0$; a square-root closure has unbounded
slope there. \textbf{This deviation is what makes the boiling front form at all} in a
differentiable formulation.

\subsubsection{D-TH-1 and D-FLOW-1}
Eulerian mixture void rather than Lagrangian slug tracking, and prescribed flow rather than
prescribed pump head.

Registered but \emph{off by default}, so no result here depends on them: condensation
(D-TH-4), three-group decay heat (D-KIN-3), axial fuel expansion (D-FB-4). Vapour expansion
accelerating the flow (D-TH-2) is implemented but not usable. Five further feedback
mechanisms are omitted and recorded as such (D-FB-3), and there is no pin failure or
material relocation (D-SCOPE-1)---which bounds the transient to its pre-failure phase.

\section{Reference solution and its resolution}

The reference is a stiff Radau integration on an axial mesh, and it is the instrument every
claim below is measured against, so its own resolution bounds what can be claimed.

It reproduces the analytic steady state to $3.4\times10^{-11}$\,K\,s$^{-1}$ and conserves
energy at second order in the step. Refining it against itself gives table~\ref{tab:ruler}.

\begin{table}[ht]
\caption{The reference solution's own error at the scoring mesh. These bound what any
comparison against it can claim.}
\label{tab:ruler}
\centering
\begin{tabular}{@{}ll@{}}
\toprule
Quantity & Reference's own error \\
\midrule
Temperatures, relative $L_2$ & 1.1--$1.6\times10^{-3}$ \\
Boiling onset time           & 0.009\,s \\
Onset height                 & 0.06 cells \\
Peak voided length           & 0.57\% \\
Pointwise void fraction      & $3.2\times10^{-2}$ \\
\bottomrule
\end{tabular}
\end{table}

Two consequences are load-bearing. First, \textbf{an acceptance bar on the pointwise void
fraction is not supportable} and was withdrawn: the reference does not know that quantity
to better than 3\%. Second, calibration practice requires a tolerance to sit at least four
times above the uncertainty of the instrument measuring it [6], so the 1\% temperature bar
is sound at a ratio of about 6 and the 0.5\,s onset criterion at about 56---but a
\emph{result} below that uncertainty is measuring the instrument.

The reference terminates when any temperature reaches the top of the section~12.13 property
fits, which defines the validity horizon; the surrogate is trained only over that horizon.

\begin{figure}[ht]
\centering
\includegraphics[width=\linewidth]{boundary_conditions.png}
\caption{What drives the transient. The pumps coast down to the natural-circulation
floor while the power is still near nominal, so the outlet temperature rises until it
meets saturation plus superheat. The dashed line marks boiling onset.}
\label{fig:bc}
\end{figure}

\textbf{Onset time is located by root-finding, not by grid inspection.} The reference's own
onset is \textbf{10.9784\,s}; read off the 0.25\,s output grid it appears as 10.75\,s. A
quarter of a second of every onset error previously reported for this model was that
quantisation.

\section{The surrogate}

Described here only as far as reproduction requires; the design study is in the repository
documentation.

The state is written multiplicatively,
\begin{equation}
\theta(\zeta,\hat{t}) = \theta_0(\zeta)\exp\bigl(\hat{t}\,N(\zeta,\hat{t})\bigr),
\end{equation}
with $\theta_0$ the analytic steady profile. This makes the initial condition exact, keeps
every temperature at or above the inlet, and pins the single upstream boundary condition
the advection equation admits---all without penalty terms. An additive ansatz, tried first,
allowed the optimiser to drive the fuel temperature negative, at which point the
logarithmic Doppler term is undefined.

Inputs pass through a random Fourier embedding [4]
$x \mapsto [\sin(2\pi Bx), \cos(2\pi Bx)]$ with $B$ frozen, at \textbf{64 features}. The
trunk is a 64-wide, 5-layer $\tanh$ MLP with five outputs, carrying \textbf{17\,029 fitted
parameters}; the embedding contributes a further 8192 in the read-out layer, whose width
follows the embedding rather than anything the network can represent, and 128 frozen
entries in $B$.

Training minimises the squared PDE residuals [2] on \textbf{one fixed set of 5000
collocation points}, with \textbf{50\,000 L-BFGS iterations} [5] under a strong-Wolfe line
search and \textbf{no first-order stage at all}. Each of those three choices is measured
rather than inherited: a first-order warm start degrades the result rather than helping it,
a set redrawn during the solve costs a factor of 1.5 against a fixed one, and 5000 points
is where the onset error falls below the reference's own resolution. At 2000 points a front
still forms but is badly placed: onset is 0.19\,s late and the saturation margin is
$+28$\,K against the reference's $+69$\,K. All arithmetic is float64 on CPU; curvature pairs are meaningless at
float32 residual magnitudes.

Two independent implementations exist, in PyTorch and JAX/Equinox, required by test to
expose identical knobs and defaults and sharing only the numpy definition of the physics.
Results below are from the JAX implementation; the cross-implementation agreement and its
limits are documented in the repository.

\section{What the surrogate reproduces}

All results at three seeds unless stated, against the reference at its scoring mesh.

\subsection{Temperature fields}

The relative $L_2$ error on the film temperature is $1.6\times10^{-3}$ (range
1.62--$1.65\times10^{-3}$), against a 1\% acceptance bar. Because the reference's own error
is 1.1--$1.6\times10^{-3}$, the surrogate has reached the resolution of the instrument
judging it: the bar is met, and \emph{how far inside} it is not a question this reference
can answer. The other fields follow at $1.7\times10^{-3}$ (coolant), $2.6\times10^{-3}$
(fuel) and $3.7\times10^{-3}$ (cladding).

\subsection{The boiling front}

\begin{figure}[ht]
\centering
\includegraphics[width=\linewidth]{temperature_map.png}
\caption{Coolant temperature over space and time. The cyan contour is
$T_{sat}+\Delta T_{sup}$, which under D-TH-3 \emph{is} the boiling front rather than a
separately computed curve; the star marks onset, located by tangency.}
\label{fig:tmap}
\end{figure}

\begin{figure}[ht]
\centering
\includegraphics[width=\linewidth]{vapor_fraction.png}
\caption{Void fraction over space and time, with axial profiles at five instants and
at onset. The front forms at the outlet and propagates downward as the coolant heats.}
\label{fig:void}
\end{figure}

\begin{figure}[ht]
\centering
\includegraphics[width=\linewidth]{front_height.png}
\caption{Front height and voided length. The saturation level set and the
$\alpha > 0.5$ contour do not coincide: the gap between them is the partially voided
region the worth integral is most sensitive to.}
\label{fig:front}
\end{figure}

The engineering quantities are reproduced (table~\ref{tab:front}).

\begin{table}[ht]
\caption{Safety-relevant front quantities against the reference.}
\label{tab:front}
\centering
\begin{tabular}{@{}lll@{}}
\toprule
Quantity & Surrogate & Reference \\
\midrule
Peak voided length           & 99.4\% of reference & --- \\
Worst-seed saturation margin & $+67.4$\,K          & $+69.2$\,K \\
Boiling onset time           & 0.0042 / 0.0046 / 0.0082\,s error & 10.9784\,s \\
\bottomrule
\end{tabular}
\end{table}

The onset criterion of 0.5\,s is met on every seed, with the worst at 0.008\,s.
\textbf{Every seed now sits below the reference's own onset uncertainty of 0.009\,s}, so
onset is reported as met and is no longer separable from the instrument either; the seed
spread has fallen to $2.0\times$ from the $32\times$ of the earlier configuration. Met is
what the measurement supports.

\textbf{Onset height is not a discriminating quantity in this model.} The coolant heats
monotonically up the channel, so the hottest point---and therefore where boiling
begins---is always the outlet, for the surrogate and the reference alike. A height
criterion here measures the mesh, not the model. An earlier revision of this work reported
that quantity as solved exactly; it had merely been restated as a tautology.

\section{What is not yet solved: the closed void-reactivity loop}

With the thermal fields prescribed, the surrogate is accurate. Driving the \emph{kinetics}
from the learned fields is a different matter, and it fails in a specific and instructive
way.

The figures in this section were measured on the earlier default configuration (256 Fourier
features, a short Adam stage, 30\,000 quasi-Newton iterations on 6000 points) and have not
been re-measured on the configuration of section~4. The cancellation ratio is a property of
the worth distribution and does not depend on the surrogate; the recovered fractions may
move.

Over the same fields and the same network, the Doppler integral is reproduced to a factor
of \textbf{1.017}. The void integral recovers only \textbf{8--16\%} of the reference's
value. The two differ in one respect: the Doppler weight has one sign, and the void worth
changes sign near the top of the core.

Splitting the void functional of~(\ref{eq:void}) at the sign change, at the instant it
peaks, gives table~\ref{tab:cancel}.

\begin{table}[ht]
\caption{The void functional split at the sign change of the worth, at peak. The two halves
are of comparable size and opposite sign.}
\label{tab:cancel}
\centering
\begin{tabular}{@{}ll@{}}
\toprule
Contribution & Value \\
\midrule
Positive-worth region, $J^{+}$ & $+4.656\times10^{-4}$ \\
Negative-worth region, $J^{-}$ & $-1.695\times10^{-4}$ \\
Their sum, $J$                 & $+2.962\times10^{-4}$ \\
Cancellation ratio, $|J|/(|J^{+}|+|J^{-}|)$ & \textbf{0.466} \\
\bottomrule
\end{tabular}
\end{table}

\begin{figure}[ht]
\centering
\includegraphics[width=0.78\linewidth]{void_worth_split.png}
\caption{The void worth $w(\zeta)$, shaded by sign. Positive over most of the core and
negative near the top, so the reactivity functional of~(\ref{eq:void}) is a difference
of two large contributions rather than a sum of small ones.}
\label{fig:worth}
\end{figure}

\begin{figure}[ht]
\centering
\includegraphics[width=0.78\linewidth]{reactivity.png}
\caption{Reactivity split by mechanism, in units of $\beta_{eff}$. Reporting only the
net would hide which component carries the error.}
\label{fig:rho}
\end{figure}

A relative error $\epsilon$ on each half therefore becomes $2.1\epsilon$ on the sum. The
functional is an ill-conditioned target by construction, and reporting it as one number
hides which half is wrong; we report the halves separately from here on.

\textbf{It is not a sampling problem.} Dual-weighted-residual theory [3] gives the error in
a functional as $\langle R(u_\theta), z^{*}\rangle$ to leading order, with $z^{*}$ the
solution of the adjoint problem sourced by the functional's derivative. For this advective
coolant operator the adjoint runs backwards in $\zeta$, and it can be evaluated in closed
form. The result is a \textbf{step}: the void slope underflows to exactly zero wherever the
coolant is subcooled, so $z^{*}$ is constant over the lower 72\% of the channel and zero
above it. Every point below the front carries equal sensitivity and every point above it
carries none. Residual-magnitude sampling concentrates points \emph{at} the
front---precisely where the functional is insensitive---so it cannot help, and a uniform
sampler is already near optimal.

Evaluated open-loop on the surrogate's own field, the functional is in fact accurate: the
positive half to 1.66--2.66\% and the negative half exactly, the latter because the
negative-worth region is fully voided in surrogate and reference alike and the integral is
then fixed by geometry. \textbf{The deficit therefore arises in the closed loop}, where
$\rho_{\mathrm{void}}$ feeds back into the kinetics and the error compounds---not in the
evaluation of the functional on a given field. That localises the remaining work precisely.

\section{Discussion}

The surrogate reproduces the thermal-hydraulic phase of this transient to the resolution of
the reference solver, including the safety-relevant front quantities: when boiling starts,
how far it extends, and by what margin. For applications that consume those
fields---uncertainty propagation over thermophysical parameters, inference of boundary
conditions from limited instrumentation, gradient-based studies of the coastdown---the
model is usable now, and its derivatives with respect to inputs are exact by construction.

For applications that require the \textbf{feedback loop closed}, it is not. The obstruction
is not accuracy in any ordinary sense---the fields are right and the open-loop functional
is right---but the conditioning of a near-cancelling integral inside a feedback path. That
is a property of the physics: a positive void worth over most of the core against a
negative worth near the top is precisely what makes SFR void feedback interesting, and it
is precisely what makes it a hard target for a surrogate. Any neural approach to sodium
void feedback will meet the same $2.1\times$ amplification.

Two consequences for practice. First, \textbf{a surrogate should be qualified against the
functionals it will be used for, not only against field norms}: an $L_2$ of
$1.6\times10^{-3}$ on temperature coexists here with an 84--92\% miss on a reactivity
integral built from the same field. Second, \textbf{the reference's own resolution must be
quantified before an acceptance bar is set} [7, 8]. Two bars in this work were set without
that check; one was withdrawn, and the temperature result has since reached the point where
the instrument, not the model, is the limit.

\section{Conclusions}

A physics-informed neural surrogate for the sodium-boiling phase of an SFR unprotected
loss-of-flow transient, built on SAS4A/SASSYS-1 thermophysics and feedback laws with four
registered deviations, reproduces the reference solution's temperature fields to
$1.6\times10^{-3}$ relative $L_2$, its boiling onset to within 0.008\,s of 10.9784\,s, and
99.4\% of its peak voided length with a saturation margin of $+67.4$\,K against
$+69.2$\,K. The first two are at the reference's own resolution and are reported as met
rather than as measured accuracies.

The closed void-reactivity loop is not reproduced, recovering 8--16\% of the reference
integral while the non-cancelling Doppler integral over the same fields is correct to
1.017. The cause is quantified---a cancellation ratio of 0.466, amplifying any error on
either half by $2.1\times$---and shown by an adjoint argument not to be a sampling
deficiency. Closing that loop is the outstanding problem for this model, and we expect it
to be the outstanding problem for neural surrogates of sodium void feedback generally.

\section*{Acknowledgments}

The SAS4A/SASSYS-1 documentation maintained by Argonne National Laboratory was the sole
source for the thermophysics and feedback laws used here.

% --- IOP numeric reference list: [n] Authors Year Title Journal Volume Page. 10 pt,
% --- 8 mm hanging indent, matching `iop-jpcs.typ`'s `references`.
\section*{References}
\begingroup
\fontsize{10}{12}\selectfont
\begin{list}{}{%
  \setlength{\leftmargin}{8mm}\setlength{\labelwidth}{8mm}
  \setlength{\itemindent}{-8mm}\setlength{\listparindent}{0mm}
  \setlength{\itemsep}{0.5em}\setlength{\parsep}{0pt}\setlength{\topsep}{0.5em}}

\item[] [1] Fanning T H (ed) 2017 \emph{The SAS4A/SASSYS-1 Safety Analysis Code System}
ANL/NE-16/19 (Argonne, IL: Argonne National Laboratory)

\item[] [2] Raissi M, Perdikaris P and Karniadakis G E 2019 Physics-informed neural
networks: a deep learning framework for solving forward and inverse problems involving
nonlinear partial differential equations \emph{J. Comput. Phys.} \textbf{378} 686

\item[] [3] Becker R and Rannacher R 2001 An optimal control approach to a posteriori error
estimation in finite element methods \emph{Acta Numerica} \textbf{10} 1

\item[] [4] Tancik M \emph{et al} 2020 Fourier features let networks learn high frequency
functions in low dimensional domains \emph{Adv. Neural Inf. Process. Syst.} \textbf{33}
7537

\item[] [5] Nocedal J and Wright S J 2006 \emph{Numerical Optimization} 2nd edn (New York:
Springer)

\item[] [6] Taylor B N and Kuyatt C E 1994 \emph{Guidelines for Evaluating and Expressing
the Uncertainty of NIST Measurement Results} NIST Technical Note 1297

\item[] [7] Jakeman J D, Barba L A, Martins J R R A and O'Leary-Roseberry T 2026
Verification and validation for trustworthy scientific machine learning
\emph{Mach. Learn.: Sci. Technol.} \textbf{7} 025055

\item[] [8] E\c{c}a L and Hoekstra M 2014 A procedure for the estimation of the numerical
uncertainty of CFD calculations based on grid refinement studies
\emph{J. Comput. Phys.} \textbf{262} 104

\end{list}
\endgroup

\end{document}
